% Generated by analysis/flagship_corrected.py -- do not edit by hand.
% Table order is cited as S33-S37 from Paper.tex; keep it stable.

\begin{table}[!t]
\centering
\caption{Agreement on the $1\%$-constrained winner, including the cross-stack comparisons that the main paper omits. $\kappa$ is Cohen's chance-corrected agreement, which de-inflates raw agreement where one mode dominates. The tie-tolerant columns count two devices as agreeing when their within-tolerance winner sets intersect; tolerance makes the within-family (S25/S26) cut exact but does not rescue any cross-stack comparison. Rows marked $\dagger$ cross the two measurement stacks and are therefore \emph{descriptive only}: a same-device stack control (Table~\ref{tab:supp-stack-control}) shows that changing the stack alone costs as much agreement as these rows report, so they must not be read as a vendor effect. All columns, including the exact one, are restricted to the modes measured and feasible on \emph{both} devices, so the exact column need not equal the per-device agreement reported in the main paper.}
\label{tab:supp-flagship-tol}
\footnotesize
\setlength{\tabcolsep}{3pt}
\begin{tabular}{@{}lrrrrrr@{}}
\toprule
Comparison & $N$ & $\kappa$ & exact & $\tau_{\mathrm{tie}}{=}2\%$ & $5\%$ & $10\%$ \\
\midrule
S25 vs S26 & 150 & 0.589 & 80.0 & 100.0 & 100.0 & 100.0 \\
S25 vs Tab~S11$^\dagger$ & 40 & 0.061 & 22.5 & 35.0 & 40.0 & 47.5 \\
S25 vs Pixel~4$^\dagger$ & 150 & 0.017 & 20.7 & 40.7 & 55.3 & 68.7 \\
Pixel~4 vs Tab~S7+ & 150 & 0.116 & 38.7 & 77.3 & 93.3 & 97.3 \\
\midrule
S25+S26+Tab~S11 agree$^\dagger$ & 40 & --- & 12.5 & --- & --- & --- \\
\bottomrule
\end{tabular}
\end{table}

\begin{table}[!t]
\centering
\caption{Regret of shipping one device's $1\%$-constrained winner on another, in relative (\%) and absolute (ms) terms. Within the Snapdragon family the penalty is negligible; across the Snapdragon--MediaTek cut it is large but, being cross-stack, is not separable from the measurement stack (Table~\ref{tab:supp-stack-control}).}
\label{tab:supp-flagship-regret}
\footnotesize
\setlength{\tabcolsep}{4pt}
\begin{tabular}{@{}lrrrrr@{}}
\toprule
Ship $\rightarrow$ target & $N$ & Med.\ \% & Mean \% & P90 \% & Med.\ ms \\
\midrule
S25 $\rightarrow$ S26 & 150 & 0.0 & 0.2 & 0.7 & 0.000 \\
S26 $\rightarrow$ S25 & 150 & 0.0 & 0.3 & 1.1 & 0.000 \\
S25 $\rightarrow$ Tab~S11 & 40 & 17.8 & 40.6 & 118.1 & 0.478 \\
Tab~S11 $\rightarrow$ S25 & 38 & 135.6 & 525.5 & 1660.9 & 0.160 \\
USB Pixel~4 $\rightarrow$ S25 & 147 & 122.4 & 485.0 & 1270.1 & 0.196 \\
USB Pixel~4 $\rightarrow$ Tab~S7+ & 150 & 0.5 & 5.1 & 23.5 & 0.031 \\
\bottomrule
\end{tabular}
\end{table}

\begin{table}[!t]
\centering
\caption{Measurement-stack control. The paper's USB harness and Qualcomm AI~Hub benchmark the same hardware on a 20-model subset each ($18$ scorable). The Pixel~4 row is a strict same-device control; on the tablet row a Galaxy Tab~S7 (SD~865) stood in on AI Hub for the paper's Tab~S7$+$ (SD~865$+$), so it is a closest-sibling rather than same-device comparison and the Pixel~4 figure is the one quoted as the stack bound. Winner-mode agreement across stacks is as low as the cross-vendor agreement of Table~\ref{tab:supp-flagship-tol}, whereas absolute latency is close to unbiased (median USB/AI~Hub ratio near unity): mode \emph{rankings}, not milliseconds, are what fail to transfer.}
\label{tab:supp-stack-control}
\footnotesize
\setlength{\tabcolsep}{3.5pt}
\begin{tabular}{@{}llrrrr@{}}
\toprule
USB device & AI Hub device & $N$ & Match \% & ms ratio & $\rho$ \\
\midrule
Galaxy Tab S7$+$ & Samsung Galaxy Tab S7$^\ddagger$ & 18 & 22.2 & 0.95 & 0.50 \\
Pixel 4 & Google Pixel 4 & 18 & 16.7 & 1.04 & 0.45 \\
\midrule
Pooled & --- & 36 & 19.4 & 0.96 & 0.50 \\
\bottomrule
\addlinespace[2pt]
\multicolumn{6}{@{}l@{}}{\footnotesize $^\ddagger$not the same physical device.} \\
\end{tabular}
\end{table}

\begin{table}[!t]
\centering
\caption{MediaTek Dimensity~9400 delegate coverage: number of the 40 subset models that produced a valid measurement on each backend, by PTQ mode. The NPU path accepts the calibrated integer graphs for almost no model, so on this device integer modes execute on the CPU or GPU. This is the mechanism behind the cross-vendor divergence in winner modes.}
\label{tab:supp-mtk-coverage}
\footnotesize
\begin{tabular}{@{}lrrrr@{}}
\toprule
Mode & CPU & GPU & NPU & Any \\
\midrule
FP32 & 40 & 39 & 34 & 40 \\
FP16 & 40 & 39 & 34 & 40 \\
dyn. & 40 & 39 & 31 & 40 \\
INT8 & 38 & 38 & 1 & 38 \\
mixed & 38 & 38 & 1 & 38 \\
\bottomrule
\end{tabular}
\end{table}


\begin{table}[!t]
\centering
\caption{Winner share (\%) per PTQ mode on the freeze-150 subset, selecting by pure latency (lat.) and under the $1\%$ accuracy budget. Dynamic INT8, the most frequent optimum on both full corpora, never wins on the AI~Hub devices; INT8's pure-latency dominance there is more than halved by the accuracy budget.}
\label{tab:supp-flagship-winrates-full}
\footnotesize
\setlength{\tabcolsep}{4pt}
\begin{tabular}{@{}llrrrrr@{}}
\toprule
Device & Sel. & FP32 & FP16 & dyn. & INT8 & mixed \\
\midrule
S25 (8 Elite) & lat. & 4.0 & 0.7 & 0.0 & 91.3 & 4.0 \\
 & $1\%$ & 34.7 & 14.0 & 0.0 & 48.7 & 2.7 \\
\addlinespace[2pt]
S26 (8 Elite Gen5) & lat. & 4.0 & 1.3 & 0.0 & 93.3 & 1.3 \\
 & $1\%$ & 32.0 & 16.7 & 0.0 & 51.3 & 0.0 \\
\addlinespace[2pt]
Tab~S11 (D9400) & lat. & 12.5 & 42.5 & 32.5 & 7.5 & 5.0 \\
 & $1\%$ & 12.5 & 45.0 & 35.0 & 5.0 & 2.5 \\
\addlinespace[2pt]
USB Pixel~4 & lat. & 25.3 & 30.0 & 25.3 & 9.3 & 10.0 \\
 & $1\%$ & 26.0 & 31.3 & 35.3 & 2.0 & 5.3 \\
\addlinespace[2pt]
USB Tab~S7+ & lat. & 22.7 & 38.7 & 32.0 & 4.7 & 2.0 \\
 & $1\%$ & 24.7 & 38.7 & 34.0 & 2.0 & 0.7 \\
\addlinespace[2pt]
USB Redmi~13 & lat. & 24.7 & 17.3 & 18.7 & 32.7 & 6.7 \\
 & $1\%$ & 26.0 & 20.0 & 33.3 & 18.0 & 2.7 \\
\addlinespace[2pt]
\bottomrule
\end{tabular}
\end{table}
